\section{Conclusion and Phase 2 Plan}\label{sec:conclusion}
In simulation on five benchmark LV networks, a transformer-side energy balance with abstention
separates power outages from communication failures for silent groups above its measured resolution
limit, and it abstains below that limit instead of guessing. Every verdict carries its evidence, a
probability, the alternative cause and a tamper-evident record. The evaluations also show where the
method fails and what deployment would cost in operator time.

\textbf{Phase 2 plan, in priority order.}
\begin{enumerate}[leftmargin=1.4em]
\item Bring operator load to at most 1.93 human items per transformer-day, so that one operator around
the clock covers a district of 149 substations (the escalating gate gives 4.15); wire the escalating
gate and the communication re-check into the escalation step; operator trials of the console.
\item Run the registered people study on the explanations.
\item Fix the evidence accounting for late stored readings and sub-interval onsets; add a gap model
that depends on load.
\item Single-meter faults: per-meter residual drift and neighbour comparison with field voltage data.
\item Shared learning through a bulletin of frozen thresholds and anonymous failure signatures.
\item Ask STEG the prerequisite questions of Section~\ref{sec:business} and run a retrospective test on
one substation.
\end{enumerate}

\textbf{Repository.} The submitted repository holds the code (MIT), the simulated dataset (ODbL 1.0,
derived from SimBench, which is ODbL 1.0), the claims ledger, the protocols and the extended report
(\texttt{docs/EXTENDED\_REPORT.md}). Quickstart: \texttt{pip install -r requirements.lock}, \texttt{python -m pytest -q},
\texttt{python -m shahid.demo}. Figures and tables contain information from the SimBench database,
made available under the ODbL 1.0.
